% !TEX root = ../main.tex
% !TEX program = lualatex
\documentclass[../main.tex]{subfiles}
\IfSubfilesClassLoaded{\externaldocument{\subfix{../../build/main}}}{}
% =============================================================================
% APPENDIX: COMBATANT COMMANDS
% DESCRIPTION: Geographic and functional combatant commands (CCMDs) reference.
% =============================================================================
\begin{document}
\IfSubfilesClassLoaded{\chapter{EDO Study Guide}}{}
%====================
% Contents here
%====================
\section{Combatant Commands (Reference)}
\noindent\textbf{\ac{aor}:} The geographic (or astrographic) area assigned to a \ac{ccdr} for missions and forces. (Title~10) % cite earlier laws in main refs if desired

\textbf{Geographic (7):}
\begin{itemize}
  \item USSPACECOM\footnote{USSPACECOM is treated by \ac{dow} as a \textbf{geographic (astrographic)} command with an AOR beginning at the K\'{a}rm\'{a}n line ($\sim$100 km)}\\
\end{itemize}
\textbf{Functional (4):}
\begin{itemize}
\item \ac{usafricom}
\item \ac{uscentcom}
\item \ac{useucom}
\item \ac{usindopacom}
\item \ac{usnorthcom}
\item \ac{ussouthcom}
\end{itemize}
%====================
% End of file
%====================
\IfSubfilesClassLoaded{
  \RenewDocumentCommand{\entryname}{}{\textbf{\color{Modern} Acronym}}
  \RenewDocumentCommand{\descriptionname}{}{\textbf{\color{Modern} Definition}}
  \printnoidxglossary[
    type=\acronymtype,
    title=Acronyms,
    style=long-booktabs
  ]}{}
\end{document}
\item \ac{ussocom}
\item \ac{usstratcom}
\item \ac{ustranscom}
\item \ac{uscybercom}
